\documentclass[10pt,a4paper]{amsart}
\usepackage[margin=19mm]{geometry}
\usepackage{amsmath,amssymb,mathtools}
\usepackage[scr=rsfs,cal=euler]{mathalfa}
\usepackage{xfrac}
\usepackage[unicode,hidelinks,bookmarks=false]{hyperref}
\newcommand{\ie}{i.e.\ }
\newcommand{\cf}{cf.\ }
\newcommand{\resp}{respectively\ }
\newcommand{\Subsection}{\subsection}
\providecommand{\nobreakdash}{\nobreak\hbox{-}}
\setlength{\emergencystretch}{2em}
\multlinegap0pt
\usepackage{bm}
\usepackage{bbm}
\usepackage{dsfont}
\usepackage{xspace}
\usepackage{cancel}
\usepackage{mathtools}
\usepackage{amsthm}
\makeatletter
\@ifundefined{theorem}{\newtheorem{theorem}{Theorem}}{}
\@ifundefined{lemma}{\newtheorem{lemma}[theorem]{Lemma}}{}
\makeatother
\def \La {\Lambda}
\def \de {\delta}
\def \la {\lambda}
\def\w {\omega}
\title[Decomposition of the Curvature Operator and Applications to ]{Decomposition of the Curvature Operator and Applications to the Hopf Conjecture}
\author{Teng Huang, Weiwei Wang}
\date{Source: arXiv:2507.15237v2 (CC BY 4.0)}
\begin{document}
\maketitle

\noindent\emph{Source and licence.} Decomposition of the Curvature Operator and Applications to the Hopf Conjecture. Version arXiv:2507.15237v2 (CC BY 4.0), distributed under the Creative Commons Attribution 4.0 International licence, as recorded in the arXiv licence metadata for this exact version and shown on the abstract page https://arxiv.org/abs/2507.15237v2. Full rights notice: licenses/arxiv-2507.15237-CC-BY-4.0.txt.\par\smallskip

\noindent\textit{Editorial scope.} Counted unit: the Lemma whose source label is \texttt{P5} in the article (source0.tex lines 332--339, source label \texttt{P5}), delivered together with the article's own complete original proof of it (source0.tex lines 340--376, one \texttt{proof} environment of 37 lines). No mathematics is written, repaired or re-derived: statement and proof are the article's own text line for line. Required context retained uncounted: none. Every definition, display and label of those lines is retained unchanged. Omitted from this package and not redistributed: 1--331, 377--1262. Nothing inside a retained excerpt is omitted or altered: every retained body is delivered exactly as the article's lines are written, so no omission receipt of any kind accompanies this package. Citations: the retained excerpts carry no citation command at all, so no bibliography entry of the article is transcribed and no reference is redistributed. Reference adaptation: none was needed and none was made; every reference inside the delivered unit resolves inside it, so no unresolved reference or citation is produced. Printed slips kept literal: no printed word, symbol, hypothesis or number of the counted statement or of its proof was altered. Layout and class: the article's own class is \texttt{article} with packages including amsmath, amsthm, amssymb, times, enumerate, bm, xy, mathrsfs, amscd, color, comment; the wrapper is the lane's pinned \texttt{amsart} editorial wrapper at 10pt on A4, which restates the theorem-like environments the retained text needs and adds the article's own macro definitions that the retained text needs (\texttt{\textbackslash La}, \texttt{\textbackslash de}, \texttt{\textbackslash la}, \texttt{\textbackslash w}), with extra wrapper packages (bm, bbm, dsfont, xspace, cancel, mathtools). The delivered numbering is the unit's own. Assets: no retained span contains a figure environment, a graphics-inclusion command, a PDF-page inclusion, a table or a TikZ picture; no image file is copied into this package, the package declares no asset dependency and no image file is redistributed with it. Licence: version arXiv:2507.15237v2 of this article is distributed under the Creative Commons Attribution 4.0 International licence, as recorded in the arXiv licence metadata for this exact version and shown on the abstract page https://arxiv.org/abs/2507.15237v2. A full rights notice is delivered at licenses/arxiv-2507.15237-CC-BY-4.0.txt. Characters: every retained excerpt is pure ASCII, so the delivered engine, XeLaTeX with its default Latin Modern text font, reports no missing character.
	\begin{lemma}\label{P5}
		Let $(M,g)$ be a compact manifold of dimension $n$, $n\geq4$, and let $\{e_{i}\}_{1,\cdots,n}$, be an orthonormal eigenbasis of the Ricci curvature with corresponding eigenvalues $\{\la_{i} \}_{1,\cdots,n}$. Then the curvature operator $\mathfrak{S}:\La^{2}TM\rightarrow\La^{2}TM$ corresponding to $S_{g}\odot g$ acts on $2$-forms as
		\begin{equation*}
		\mathfrak{S}(e_{i}\wedge e_{j})=\frac{1}{n-2}\left(\la_{i}+\la_{j}-\frac{R_{g}}{n-1}\right)(e_{i}\wedge e_{j}).
		\end{equation*}
		It implies that $\{e_{i}\wedge e_{j}\}_{1\leq i<j\leq n}$ is an  an eigenbasis of $\mathfrak{S}$ with eigenvalues
		$\la_{ij}:=\frac{1}{n-2}(\la_{i}+\la_{j}-\frac{R_{g}}{n-1})$.
	\end{lemma}

	\begin{proof}
		The $\{e_{i}\}_{i=1,\cdots,n}$ is an orthonormal eigenbasis of Ricci curvature tensor $Ric_{g}$, which satisfies
		\begin{equation*}
		Ric_{g}(e_{i})=\la_{i}e_{i},\ \quad g(e_{i},e_{j})=\de_{ij}.
		\end{equation*}
		By the definition of  Kulkarni-Nomizu product, we compute
		\begin{equation}\label{Ricg1}
		\begin{split}
		(Ric\odot g)(e_{i},e_{j},e_{k},e_{l})&=Ric_{g}(e_{i},e_{k})g(e_{j},e_{l})-Ric_{g}(e_{i},e_{l})g(e_{j},e_{k})\\
		&+Ric_{g}(e_{j},e_{l})g(e_{i},e_{k})-Ric_{g}(e_{j},e_{k})g(e_{i},e_{l})\\
		&=\la_{i}\de_{ik}\de_{jl}-\la_{i}\de_{il}\de_{jk}+\la_{j}\de_{jl}\de_{ik}-\la_{j}\de_{jk}\de_{il}\\
		&=(\la_{i}+\la_{j})(\de_{ik}\de_{jl}-\de_{il}\de_{jk}).\\
		\end{split}
		\end{equation}
		Substituting \eqref{Ricg1} into $S_{g}\odot g$, we obtain
		\begin{equation*}
		\begin{split}
		(S_{g}\odot g)(e_{i},e_{j},e_{k},e_{l})&=\left[\frac{1}{n-2}\left(Ric_{g}-\frac{R_{g}}{2(n-1)}\cdot g\right)\odot g\right](e_{i},e_{j},e_{k},e_{l})\\
		&=\frac{1}{n-2}\left(\la_{i}+\la_{j}-\frac{R_{g}}{n-1}\right)(\de_{ik}\de_{jl}-\de_{il}\de_{jk}).
		\end{split}
		\end{equation*}	
		For the curvature operator $\mathfrak{S}:\La^{2}TM\rightarrow\La^{2}TM$ associated to $S_{g} \odot g$ and $\w=\sum_{i,j}\w_{ij}e_i\wedge e_j$,  the $e_{k}\wedge e_{l}(k<l)$ component of $\mathfrak{S}(\w)$  is
		$$ (\mathfrak{S}(\w))_{kl}=\sum_{i,j}(S_{g}\odot g)(e_{i},e_{j},e_{k},e_{l})\w_{ij}.$$
		Thus
		$$\mathfrak{S}(e_{i}\wedge e_{j})_{kl}=\frac{1}{n-2}(\la_{i}+\la_{j}-\frac{R_{g}}{n-1})(\de_{ik}\de_{jl}-\de_{il}\de_{jk}).$$
		\begin{itemize}
			\item if $(i,j)=(k,l)$, then  
			$$\mathfrak{S}(e_{i}\wedge e_{j})_{ij}=\frac{1}{n-2}(\la_{i}+\la_{j}-\frac{R_{g}}{n-1}).$$
			\item if $(i,j)\neq (k,l)$, then
			$$\mathfrak{S}(e_{i}\wedge e_{j})_{kl}=0.$$
		\end{itemize}
		Since  
		$\{e_{i}\wedge e_{j}\}_{1\leq i<j\leq n}$ forms an orthonormal basis for $\La^{2}TM$, it follows that this set is an eigenbasis of $\mathfrak{S}$ with eigenvalues
		$\la_{ij}:=\frac{1}{n-2}(\la_{i}+\la_{j}-\frac{R_{g}}{n-1})$, which proves the lemma.
		%is  defined by:
		%$$g(\mathfrak{S}(X\wedge Y),Z\wedge W)=(S_{g}\odot g)(X,Y,Z,W),$$
	\end{proof}

\end{document}
